
\begin{center}
\begin{minipage}{0.48\textwidth}
\includegraphics[width=\textwidth]{/home/jeff/Documents/IUTOrsay/1-Cours/4p13-Alliages/Cours-TD/Exos/Binaires-Lecture/Homogene-Ex01/figures/Fig01}
\end{minipage}\begin{minipage}{0.48\textwidth}
\captionof{figure}{Diagramme binaire isobare, à une pression de \SI{1}{\bar}, du mélange \ce{Sb}-\ce{Bi}.}
[La composition est donnée par la fraction massique en antimoine dans le mélange]
\end{minipage}
\end{center}

Le diagramme binaire isobare ($P = 1 ~ \si{\bar}$) du bismuth \ce{Bi} et de l'antimoine \ce{Sb} est donné ci-dessous, la composition étant exprimée en fraction massique de \ce{Sb}, soit $w_{\ce{Sb}}$.

\Q Un mélange solide de bismuth et d'antimoine est chauffé progressivement en milieu fermé. Il commence à fondre à $\theta = 400 ~ \si{\celsius}$. En déduire :

\begin{enumerate}
\item sa composition en fraction massique, puis en fraction molaire ;

\item la composition, en fraction massique, de la première goutte qui apparaît ;

\item la composition, en fraction massique, du dernier cristal qui disparaît et la
température de fin de fusion. 
\end{enumerate}

\Q Un mélange équimolaire en bismuth et antimoine, de masse totale $m = 100 ~ \si{\gram}$, est refroidi jusqu'à la température $\theta = 400 ~ \si{\celsius}$. Déterminer :
\begin{enumerate}
\item sa température de début d'apparition des premiers cristaux ;

\item la nature et la composition en fraction massique des phases présentes à l'équilibre ;

\item la masse des différentes phases, ainsi que la masse d'antimoine \ce{Sb} dans chacune de cès phases.
\end{enumerate}


\paragraph*{Données} 
\begin{align*}
M(\ce{Bi}) &= 209 ~ \si{\gram\per\mole} \\
M(\ce{Sb}) &= 122 ~ \si{\gram\per\mole} \\
\end{align*}